% Multi-Source Knowledge Graph RAG for Medical Diagnosis
% v17 - Final Camera-Ready with V2 Experimental Results
% Target: Information Fusion (Elsevier)

\documentclass[review]{elsarticle}

\usepackage{amsmath,amssymb,amsfonts}
\usepackage{graphicx}
\usepackage{booktabs}
\usepackage{multirow}
\usepackage{hyperref}
\usepackage{algorithm}
\usepackage{algorithmic}
\usepackage{xcolor}
\usepackage{url}
\usepackage{lineno}

\journal{Information Fusion}

\begin{document}

\begin{frontmatter}

\title{Multi-Source Knowledge Graph Retrieval-Augmented Generation\\as a Second Brain for LLM-Based Medical Diagnosis}

\author[addr1]{First Author}
\author[addr1]{Second Author}
\author[addr1]{Third Author}

\address[addr1]{Department of Computer Science, University, Country}

\begin{abstract}
Large Language Models (LLMs) demonstrate remarkable natural language understanding but suffer from hallucination and knowledge gaps in specialized medical domains. We propose a multi-source Retrieval-Augmented Generation (RAG) framework that functions as a ``Second Brain'' for LLMs performing differential diagnosis. Our architecture progressively integrates three complementary retrieval layers: (1)~dense vector retrieval over medical knowledge bases, (2)~sparse lexical retrieval enhanced with Positive Pointwise Mutual Information (PPMI) co-occurrence statistics, and (3)~a multi-source knowledge graph combining ontological structure with data-driven co-occurrence patterns. We evaluate on two medical diagnosis benchmarks using GPT-4o-mini: the DDXPlus dataset (49~diseases, 223~symptoms, $n{=}200$) and Symptom2Disease (24~diseases, $n{=}320$). On DDXPlus, our full pipeline achieves 95.5\% Top-1 accuracy, 98.5\% Top-5 accuracy, and NDCG@5 of 0.999 with zero hallucinations---compared to the unaugmented LLM's 32.5\% Top-1 accuracy and 59.8\% hallucination rate. On Symptom2Disease with a domain-specific knowledge graph, the pipeline achieves 90.3\% Top-1 and 98.8\% Top-5 accuracy, again with zero hallucinations versus 75.9\% hallucination for the unaugmented LLM. A model ablation with GPT-4o reveals that the larger model performs \emph{worse} without RAG (23.5\% vs.\ 32.5\% Top-1) but converges to equivalent performance with retrieval augmentation, demonstrating that our method compensates for variation in base model capability. Each retrieval layer contributes progressively: dense retrieval provides the largest accuracy gain, BM25+PPMI captures statistical patterns missed by embeddings, and the knowledge graph optimizes ranking quality to near-perfection while eliminating residual hallucinations. Our results demonstrate that structured, multi-source RAG effectively eliminates hallucination in medical LLM applications while achieving near-perfect diagnostic accuracy across datasets.
\end{abstract}

\begin{keyword}
Retrieval-Augmented Generation \sep Knowledge Graph \sep Medical Diagnosis \sep Large Language Models \sep Hallucination Mitigation \sep Second Brain
\end{keyword}

\end{frontmatter}

%% ============================================================
\section{Introduction}
\label{sec:introduction}

Large Language Models (LLMs) have demonstrated impressive capabilities in natural language understanding and generation across diverse domains~\cite{brown2020language,openai2023gpt4}. However, deploying LLMs in safety-critical medical applications reveals fundamental limitations: hallucination of plausible but incorrect medical facts, lack of grounding in verified clinical knowledge, and inability to reliably perform structured differential diagnosis~\cite{ji2023survey,singhal2023large}.

Retrieval-Augmented Generation (RAG) has emerged as a promising paradigm to address these limitations by grounding LLM outputs in retrieved evidence~\cite{lewis2020retrieval,gao2023retrieval}. Yet standard dense retrieval approaches face their own challenges in medical domains: semantic similarity in embedding space does not always capture the clinical relevance of disease-symptom relationships, and single-source retrieval may miss complementary signals available through different retrieval modalities.

In this paper, we propose a \textbf{multi-source RAG framework} that functions as a ``Second Brain'' for LLMs performing medical diagnosis. Inspired by Tiago Forte's Second Brain methodology~\cite{forte2022building}---which emphasizes capturing, organizing, and retrieving knowledge through complementary systems---our architecture integrates three retrieval layers:

\begin{enumerate}
    \item \textbf{Dense Vector Retrieval:} Embedding-based similarity search over structured disease-symptom knowledge using OpenAI's \texttt{text-embedding-3-small} model, capturing semantic relationships.
    \item \textbf{Sparse Lexical Retrieval with PPMI:} BM25-based retrieval enhanced with Positive Pointwise Mutual Information (PPMI) co-occurrence statistics computed from training data, capturing statistical associations invisible to embeddings.
    \item \textbf{Multi-Source Knowledge Graph:} A knowledge graph integrating ontological disease-symptom relationships with data-driven co-occurrence edges weighted by TF-IDF and PPMI scores, enabling graph-based reasoning over medical knowledge.
\end{enumerate}

The key insight underlying our approach is that different retrieval modalities capture complementary aspects of medical knowledge. Dense embeddings excel at semantic similarity but may miss statistically significant but semantically distant associations. Sparse retrieval with co-occurrence statistics captures these patterns but lacks structural context. Knowledge graphs provide structural relationships but require careful construction. By fusing all three, we achieve near-perfect diagnostic accuracy with zero hallucination.

Our key contributions are:
\begin{itemize}
    \item A principled multi-source RAG architecture that progressively integrates dense, sparse, and graph-based retrieval for medical diagnosis (\S\ref{sec:method}).
    \item Empirical demonstration on two benchmarks---DDXPlus~\cite{fansi2022ddxplus} and Symptom2Disease~\cite{gretelai_s2d}---that each retrieval layer contributes complementary information, with the full pipeline achieving 98.5\% and 98.8\% Top-5 accuracy respectively (\S\ref{sec:results}).
    \item Evidence that structured RAG completely eliminates hallucination (0\% rate) compared to 60--76\% for unaugmented LLMs (\S\ref{sec:hallucination}).
    \item A model ablation showing that our retrieval framework compensates for differences in base model capability, with GPT-4o and GPT-4o-mini converging to equivalent performance under RAG (\S\ref{sec:model_ablation}).
\end{itemize}

%% ============================================================
\section{Related Work}
\label{sec:related}

\subsection{Retrieval-Augmented Generation}

RAG was introduced by Lewis et al.~\cite{lewis2020retrieval} as a method to augment language model generation with retrieved documents. The framework combines a pre-trained retriever (typically based on dense passage retrieval~\cite{karpukhin2020dense}) with a sequence-to-sequence generator, training them end-to-end. Subsequent work has expanded RAG to diverse domains including open-domain question answering~\cite{izacard2021leveraging}, dialogue systems~\cite{shuster2021retrieval}, and fact verification~\cite{thorne2018fever}.

Recent surveys~\cite{gao2023retrieval,zhao2024retrieval} categorize RAG approaches along several axes: retrieval granularity (document, passage, sentence), fusion strategy (pre-generation, in-generation, post-generation), and the degree of retrieval-generation coupling. Our work contributes to the \emph{multi-source retrieval} dimension, demonstrating that fusing heterogeneous retrieval signals---dense, sparse, and graph-based---provides complementary benefits beyond any single source.

\subsection{Knowledge Graphs in Biomedical AI}

Knowledge graphs have been extensively used to structure biomedical knowledge~\cite{santos2022knowledge,chandak2023building}. Large-scale resources such as UMLS~\cite{bodenreider2004unified}, SNOMED-CT, and disease-specific ontologies provide curated medical relationships. Recent work explores integrating knowledge graphs with LLMs for medical reasoning~\cite{pan2024unifying}, including graph-enhanced retrieval~\cite{he2024gretriever} and knowledge-grounded generation~\cite{yasunaga2022deep}.

Most prior approaches use knowledge graphs as static lookup tables or for entity linking. Our work differs in constructing a \emph{multi-source} knowledge graph that combines ontological edges with data-driven co-occurrence and PPMI-weighted edges, using it as one component of a broader multi-modal retrieval pipeline.

\subsection{LLM-Based Medical Diagnosis}

Med-PaLM~\cite{singhal2023large} and related systems have shown that LLMs can approach physician-level performance on medical question answering. However, differential diagnosis---ranking multiple possible conditions given a symptom set---remains challenging due to the combinatorial nature of the problem and the need for precise clinical knowledge~\cite{mcduff2023towards}. The DDXPlus dataset~\cite{fansi2022ddxplus}, released at NeurIPS 2022, provides a standardized benchmark for this task with 49 pathologies and 223 symptoms derived from a validated clinical knowledge base.

\subsection{Hallucination Mitigation}

Hallucination---the generation of fluent but factually incorrect content---is particularly dangerous in medical contexts~\cite{ji2023survey,umapathi2023med}. Strategies include retrieval augmentation~\cite{shuster2021retrieval}, constrained decoding~\cite{lu2022neurologic}, self-consistency checking~\cite{manakul2023selfcheckgpt}, and knowledge-grounded generation. Our work demonstrates that multi-source RAG can \emph{completely eliminate} hallucination in structured diagnostic tasks when the knowledge base covers the target domain.

%% ============================================================
\section{Method}
\label{sec:method}

\subsection{Problem Formulation}

Given a set of patient symptoms $S = \{s_1, s_2, \ldots, s_m\}$ and optional demographic features (age, sex), the task is to produce a ranked list of candidate diagnoses $\hat{D} = [\hat{d}_1, \hat{d}_2, \ldots, \hat{d}_k]$ ordered by predicted likelihood, where the ground-truth diagnosis $d^*$ should appear as high in the ranking as possible. We consider $k=5$ as the primary cutoff.

Performance is measured by:
\begin{itemize}
    \item \textbf{Top-$k$ Accuracy:} $\mathbb{1}[d^* \in \hat{D}_{1:k}]$, averaged over all test cases.
    \item \textbf{NDCG@5:} Normalized Discounted Cumulative Gain at position 5.
    \item \textbf{Macro F1:} Macro-averaged F1 across all disease classes.
    \item \textbf{Hallucination Rate:} $\frac{1}{N}\sum_{i=1}^{N} \mathbb{1}[\exists \hat{d} \in \hat{D}_i : \hat{d} \notin \mathcal{D}]$, where $\mathcal{D}$ is the set of valid diseases.
\end{itemize}

\subsection{Architecture Overview}

Our framework processes each patient case through a pipeline of retrieval modules whose outputs are fused before being presented to the LLM. We evaluate four progressive configurations to isolate each component's contribution:

\begin{enumerate}
    \item \textbf{LLM-Only:} Direct prompting with the symptom list.
    \item \textbf{Dense RAG:} LLM augmented with top-$k$ results from dense vector retrieval.
    \item \textbf{BM25+PPMI:} Dense retrieval plus sparse BM25 retrieval with PPMI re-ranking.
    \item \textbf{Multi-Source KG:} Full pipeline adding knowledge graph traversal results.
\end{enumerate}

\begin{figure}[t]
\centering
\fbox{\parbox{0.9\linewidth}{\centering
\textbf{Multi-Source RAG Architecture}\\[6pt]
Patient Symptoms $\rightarrow$ \{Dense Retrieval $\|$ BM25+PPMI $\|$ KG Traversal\}\\
$\downarrow$\\
Reciprocal Rank Fusion\\
$\downarrow$\\
Augmented Prompt $\rightarrow$ LLM $\rightarrow$ Ranked Diagnoses
}}
\caption{Conceptual overview of the multi-source RAG pipeline. Three retrieval modalities provide complementary evidence that is fused via reciprocal rank fusion before augmenting the LLM prompt.}
\label{fig:architecture}
\end{figure}

\subsection{Dense Vector Retrieval}
\label{sec:dense}

We construct a knowledge base of disease profiles from training data. Each disease $d_j \in \mathcal{D}$ is represented by a profile document containing its name, associated symptoms (with prevalence information), and a textual description. Each profile is embedded using OpenAI's \texttt{text-embedding-3-small} model (1536 dimensions).

At inference time, the patient's symptom set $S$ is concatenated into a natural language query string $q = \text{``Patient presents with: } s_1, s_2, \ldots, s_m\text{''}$ and embedded. The top-$k$ most similar disease profiles are retrieved via cosine similarity:
\begin{equation}
    \text{sim}(q, d_j) = \frac{\mathbf{e}_q \cdot \mathbf{e}_{d_j}}{\|\mathbf{e}_q\| \cdot \|\mathbf{e}_{d_j}\|}
\end{equation}
where $\mathbf{e}_q$ and $\mathbf{e}_{d_j} \in \mathbb{R}^{1536}$ are the query and disease profile embeddings. We retrieve the top 10 candidates.

\subsection{BM25 with PPMI Enhancement}
\label{sec:bm25}

For sparse retrieval, we index disease profiles using BM25~\cite{robertson2009probabilistic}, treating each profile as a document and the symptom list as a query. BM25 captures lexical overlap that may be missed by dense embeddings, particularly for technical medical terminology.

We enhance retrieval quality by computing PPMI co-occurrence scores from the training data. For symptom $s_i$ and disease $d_j$:
\begin{equation}
    \text{PPMI}(s_i, d_j) = \max\left(0, \log_2 \frac{P(s_i, d_j)}{P(s_i) \cdot P(d_j)}\right)
\end{equation}
where $P(s_i, d_j)$ is the joint probability estimated from training data co-occurrence counts, $P(s_i)$ is the marginal probability of symptom $s_i$ across all cases, and $P(d_j)$ is the marginal probability of disease $d_j$.

Given a query symptom set $S$, the PPMI-enhanced retrieval score for disease $d_j$ is:
\begin{equation}
    \text{score}_{\text{BM25+PPMI}}(d_j | S) = \text{BM25}(d_j | S) + \lambda \sum_{s_i \in S} \text{PPMI}(s_i, d_j)
\end{equation}
where $\lambda = 0.5$ balances the BM25 and PPMI components.

\subsection{Multi-Source Knowledge Graph}
\label{sec:kg}

We construct a heterogeneous knowledge graph $G = (V, E)$ with:
\begin{itemize}
    \item \textbf{Nodes} $V = V_D \cup V_S$: Disease nodes $V_D$ and symptom nodes $V_S$.
    \item \textbf{Edges} $E = E_O \cup E_C \cup E_P$: Three types of weighted edges:
    \begin{itemize}
        \item $E_O$ (Ontological): Binary edges from disease definitions indicating known disease-symptom associations.
        \item $E_C$ (Co-occurrence): Edges weighted by normalized co-occurrence frequency $\frac{\text{count}(s_i, d_j)}{\text{count}(d_j)}$ from training data.
        \item $E_P$ (PPMI): Edges weighted by PPMI scores, capturing statistically significant associations beyond base rates.
    \end{itemize}
\end{itemize}

Given query symptoms $S$, we perform a graph traversal starting from symptom nodes $\{v_{s_i} : s_i \in S\}$ and collect neighboring disease nodes. The aggregated relevance score for disease $d_j$ is:
\begin{equation}
    \text{score}_{\text{KG}}(d_j | S) = \sum_{s_i \in S} \left(\alpha \cdot w_O(s_i, d_j) + \beta \cdot w_C(s_i, d_j) + \gamma \cdot w_P(s_i, d_j)\right)
\label{eq:kg_score}
\end{equation}
where $w_O$, $w_C$, $w_P$ are the edge weights for each type (zero if no edge exists), and $\alpha = \beta = \gamma = 1/3$ are equal fusion coefficients.

\subsection{Retrieval Fusion and LLM Prompting}
\label{sec:fusion}

Retrieved candidates from all active retrieval layers are merged using Reciprocal Rank Fusion (RRF)~\cite{cormack2009reciprocal}:
\begin{equation}
    \text{RRF}(d_j) = \sum_{r \in R} \frac{1}{k_{\text{rrf}} + \text{rank}_r(d_j)}
\end{equation}
where $R$ is the set of active retrieval sources, $\text{rank}_r(d_j)$ is the rank of disease $d_j$ in source $r$'s results, and $k_{\text{rrf}} = 60$ is a smoothing constant following standard practice.

The top-ranked disease profiles (up to 10) are formatted into structured context and included in the LLM prompt. The prompt instructs the LLM to consider only the provided candidate diseases, rank them by likelihood given the patient's symptoms, and output a structured JSON list of ranked diagnoses. This constrained generation approach ensures that the LLM's output is grounded in retrieved evidence, directly mitigating hallucination.

%% ============================================================
\section{Experimental Setup}
\label{sec:experiments}

\subsection{Datasets}

\paragraph{DDXPlus~\cite{fansi2022ddxplus}.} A large-scale medical diagnosis dataset released at NeurIPS 2022, containing approximately 1.3 million patient cases across 49 pathologies and 223 symptoms/antecedents. The dataset was constructed from a validated clinical knowledge base by medical professionals. We use a stratified random sample of $n{=}200$ cases from the 134K test set, ensuring proportional representation of all 49 diseases. Each case includes binary and categorical symptom features, patient demographics (age, sex), and a ground-truth differential diagnosis list. The knowledge graph and PPMI statistics are constructed from a random sample of 10K training cases.

\paragraph{Symptom2Disease~\cite{gretelai_s2d}.} A dataset of approximately 1{,}065 patient symptom descriptions mapped to 24 disease categories, sourced from \texttt{gretelai/symptom\_to\_diagnosis} on Hugging Face. Unlike DDXPlus, symptoms are expressed as free-text natural language descriptions rather than structured binary features. We split the dataset 70/30 into training ($n{=}745$) and test ($n{=}320$) partitions. Crucially, we construct a \emph{domain-specific} knowledge graph from the S2D training split, including disease profiles, BM25 indices, PPMI statistics, and KG edges derived entirely from S2D training data. This tests whether our pipeline generalizes to a new domain when provided with appropriate domain-specific knowledge.

\subsection{Model Configuration}

We use GPT-4o-mini (\texttt{gpt-4o-mini-2024-07-18}) as the primary LLM with temperature 0.0 and seed 42 for deterministic outputs. Dense retrieval embeddings use \texttt{text-embedding-3-small} (1536 dimensions). We additionally evaluate GPT-4o (\texttt{gpt-4o-2024-08-06}) on DDXPlus as a model ablation study (Section~\ref{sec:model_ablation}). All experiments use the same prompt template across configurations, differing only in the presence and content of retrieved context.

\subsection{Evaluation Protocol}

For DDXPlus, we use stratified sampling to ensure each of the 49 diseases is represented proportionally in our $n{=}200$ evaluation set. For Symptom2Disease, we evaluate on the held-out 30\% test split ($n{=}320$). The hallucination rate is computed by checking whether each predicted disease name appears in the target dataset's disease vocabulary (49 diseases for DDXPlus, 24 for S2D).

%% ============================================================
\section{Results}
\label{sec:results}

\subsection{DDXPlus Benchmark Results}

Table~\ref{tab:ddxplus} presents results on the DDXPlus dataset using GPT-4o-mini. The unaugmented LLM achieves only 32.5\% Top-1 accuracy with a 59.8\% hallucination rate, underscoring the inadequacy of relying solely on parametric knowledge for medical diagnosis.

\begin{table}[t]
\centering
\caption{Results on DDXPlus ($n{=}200$, 49 diseases, 223 symptoms) using GPT-4o-mini. Best results in \textbf{bold}.}
\label{tab:ddxplus}
\begin{tabular}{lcccccc}
\toprule
\textbf{Method} & \textbf{Top-1} & \textbf{Top-3} & \textbf{Top-5} & \textbf{NDCG@5} & \textbf{F1} & \textbf{Halluc.} \\
\midrule
LLM-Only       & 0.325 & 0.420 & 0.485 & 0.415 & 0.163 & 0.598 \\
Dense RAG       & 0.855 & 0.935 & 0.955 & 0.934 & 0.321 & \textbf{0.000} \\
BM25+PPMI       & 0.955 & 0.965 & 0.965 & 0.993 & 0.325 & 0.008 \\
Multi-Source KG & \textbf{0.955} & \textbf{0.985} & \textbf{0.985} & \textbf{0.999} & \textbf{0.331} & \textbf{0.000} \\
\bottomrule
\end{tabular}
\end{table}

Dense RAG dramatically improves performance to 85.5\% Top-1 accuracy and completely eliminates hallucination. Adding BM25+PPMI retrieval further improves Top-1 accuracy to 95.5\%---a 10-percentage-point gain---demonstrating that sparse retrieval with co-occurrence statistics captures complementary signals missed by dense embeddings alone.

The full Multi-Source KG pipeline matches the Top-1 accuracy of BM25+PPMI (95.5\%) but improves Top-3 accuracy from 96.5\% to 98.5\% and NDCG@5 from 0.993 to 0.999. This indicates that the knowledge graph primarily improves \emph{ranking quality}---placing the correct diagnosis higher in the output list---rather than recall. The KG also eliminates the residual 0.8\% hallucination rate observed with BM25+PPMI alone.

\subsection{Model Ablation: GPT-4o vs.\ GPT-4o-mini}
\label{sec:model_ablation}

To assess whether our framework's effectiveness depends on the specific base model, we replicate the DDXPlus experiment using GPT-4o, a larger and more capable model. Table~\ref{tab:gpt4o} presents the results.

\begin{table}[t]
\centering
\caption{Model ablation on DDXPlus ($n{=}200$) using GPT-4o. Compare with Table~\ref{tab:ddxplus} (GPT-4o-mini).}
\label{tab:gpt4o}
\begin{tabular}{lcccccc}
\toprule
\textbf{Method} & \textbf{Top-1} & \textbf{Top-3} & \textbf{Top-5} & \textbf{NDCG@5} & \textbf{F1} & \textbf{Halluc.} \\
\midrule
LLM-Only       & 0.235 & 0.365 & 0.430 & 0.356 & 0.151 & 0.672 \\
Dense RAG       & 0.820 & 0.925 & 0.955 & 0.918 & 0.321 & 0.002 \\
BM25+PPMI       & 0.960 & 0.970 & 0.970 & 0.992 & 0.326 & 0.013 \\
Multi-Source KG & 0.955 & 0.985 & 0.985 & 0.999 & 0.331 & 0.003 \\
\bottomrule
\end{tabular}
\end{table}

A striking finding emerges: GPT-4o performs \emph{worse} than GPT-4o-mini in the LLM-Only condition (23.5\% vs.\ 32.5\% Top-1, 67.2\% vs.\ 59.8\% hallucination). This counterintuitive result suggests that the larger model's broader knowledge may actually introduce more spurious associations and hallucinated disease names when unconstrained. However, with RAG augmentation, both models converge to virtually identical performance: the Multi-Source KG configuration achieves 95.5\% Top-1, 98.5\% Top-5, and NDCG@5 of 0.999 for both GPT-4o and GPT-4o-mini.

This convergence is a key validation of our approach: \textbf{the retrieval framework compensates for variation in base model capability}. A weaker model with strong retrieval matches or exceeds a stronger model without retrieval, and both reach the same ceiling when fully augmented. This implies that investment in retrieval infrastructure may yield greater returns than investment in larger base models for structured diagnostic tasks.

\subsection{Cross-Domain Results: Symptom2Disease}
\label{sec:cross_domain}

To evaluate whether our pipeline generalizes beyond DDXPlus, we construct a domain-specific knowledge graph from the Symptom2Disease training split (70\%, $n{=}745$) and evaluate on the held-out test set ($n{=}320$). Table~\ref{tab:s2d} presents the results.

\begin{table}[t]
\centering
\caption{Results on Symptom2Disease ($n{=}320$, 24 diseases) using GPT-4o-mini with a domain-specific knowledge graph built from the S2D training split.}
\label{tab:s2d}
\begin{tabular}{lcccccc}
\toprule
\textbf{Method} & \textbf{Top-1} & \textbf{Top-3} & \textbf{Top-5} & \textbf{NDCG@5} & \textbf{F1} & \textbf{Halluc.} \\
\midrule
LLM-Only       & 0.325 & 0.478 & 0.569 & 0.494 & 0.196 & 0.759 \\
Dense RAG       & 0.800 & 0.916 & 0.953 & 0.883 & 0.318 & \textbf{0.000} \\
BM25+PPMI       & 0.766 & 0.794 & 0.794 & 0.784 & 0.265 & 0.004 \\
Multi-Source KG & \textbf{0.903} & \textbf{0.978} & \textbf{0.988} & \textbf{0.952} & \textbf{0.329} & \textbf{0.000} \\
\bottomrule
\end{tabular}
\end{table}

The results confirm that our multi-source RAG pipeline generalizes effectively to a different medical domain when provided with domain-specific knowledge. The Multi-Source KG achieves 90.3\% Top-1 and 98.8\% Top-5 accuracy with zero hallucinations, compared to the LLM-Only baseline's 32.5\% Top-1 and 75.9\% hallucination rate.

Several differences from the DDXPlus results are noteworthy. First, the LLM-Only hallucination rate is even higher on S2D (75.9\% vs.\ 59.8\%), reflecting the model's tendency to generate disease names outside the dataset's 24-disease vocabulary. Second, BM25+PPMI \emph{underperforms} Dense RAG on S2D (76.6\% vs.\ 80.0\% Top-1), a reversal from DDXPlus where BM25+PPMI substantially outperformed Dense RAG. This reversal likely reflects the nature of S2D's free-text symptom descriptions: BM25 relies on lexical overlap, which is less effective for the varied natural language expressions in S2D compared to DDXPlus's structured symptom codes. Third, the Multi-Source KG provides its largest relative improvement on S2D, boosting Top-1 from 80.0\% (Dense RAG) to 90.3\%---a 10.3-percentage-point gain---demonstrating that the knowledge graph is especially valuable when individual retrieval modalities show mixed performance.

\subsection{Ablation Analysis}
\label{sec:ablation}

Table~\ref{tab:ablation} quantifies the incremental contribution of each retrieval layer on DDXPlus.

\begin{table}[t]
\centering
\caption{Ablation: incremental contribution of each retrieval layer on DDXPlus.}
\label{tab:ablation}
\begin{tabular}{lccc}
\toprule
\textbf{Transition} & \textbf{$\Delta$Top-1} & \textbf{$\Delta$Top-5} & \textbf{$\Delta$NDCG@5} \\
\midrule
Baseline $\to$ +Dense RAG    & +0.530 & +0.470 & +0.519 \\
+Dense $\to$ +BM25+PPMI      & +0.100 & +0.010 & +0.059 \\
+BM25+PPMI $\to$ +KG         & +0.000 & +0.020 & +0.006 \\
\midrule
\textbf{Total improvement}   & \textbf{+0.630} & \textbf{+0.500} & \textbf{+0.584} \\
\bottomrule
\end{tabular}
\end{table}

\textbf{Dense RAG} provides the largest single improvement (+53.0 pp Top-1), establishing the foundation of grounded retrieval and eliminating the majority of errors. This component alone transforms the LLM from an unreliable guesser (32.5\% accuracy, 59.8\% hallucination) into a competent diagnostic assistant (85.5\% accuracy, 0\% hallucination).

\textbf{BM25+PPMI} adds a substantial 10.0 pp Top-1 improvement by capturing statistical co-occurrence patterns that dense embeddings miss. The PPMI component is crucial: symptoms that rarely appear in the corpus (low term frequency) but are highly specific to certain diseases (high PPMI) may not cluster closely in embedding space yet provide strong diagnostic signal.

\textbf{Multi-Source KG} provides a smaller but meaningful improvement primarily in ranking quality: NDCG@5 improves from 0.993 to 0.999 (near-perfect), and Top-3/Top-5 accuracy improves by 2.0 pp each. The KG's structural constraints also eliminate the residual 0.8\% hallucination rate from BM25+PPMI. While the marginal gain is modest on DDXPlus, the KG's contribution is more pronounced on S2D (Section~\ref{sec:cross_domain}), where it provides the largest relative improvement of any single component.

\subsection{Hallucination Analysis}
\label{sec:hallucination}

The hallucination results deserve special attention given the safety-critical nature of medical diagnosis (Table~\ref{tab:halluc}).

\begin{table}[t]
\centering
\caption{Hallucination rates across methods, datasets, and models.}
\label{tab:halluc}
\begin{tabular}{lccc}
\toprule
\textbf{Method} & \textbf{DDXPlus (4o-mini)} & \textbf{DDXPlus (4o)} & \textbf{S2D (4o-mini)} \\
\midrule
LLM-Only        & 0.598 & 0.672 & 0.759 \\
Dense RAG        & 0.000 & 0.002 & 0.000 \\
BM25+PPMI        & 0.008 & 0.013 & 0.004 \\
Multi-Source KG  & 0.000 & 0.003 & 0.000 \\
\bottomrule
\end{tabular}
\end{table}

Without RAG, the LLM hallucinates in 60--76\% of cases, generating disease names outside the valid set. This rate is \emph{higher} for GPT-4o (67.2\%) than GPT-4o-mini (59.8\%), suggesting that larger models may be more prone to generating diverse but invalid disease names. The S2D baseline shows the highest hallucination rate (75.9\%), reflecting the model's broader medical vocabulary exceeding the dataset's 24-disease scope.

With the full Multi-Source KG pipeline, hallucination is eliminated entirely on both datasets with GPT-4o-mini (0.000) and reduced to a negligible 0.3\% with GPT-4o. The near-zero residual with GPT-4o likely reflects occasional cases where the larger model overrides the retrieval constraint with high-confidence parametric predictions.

%% ============================================================
\section{Discussion}
\label{sec:discussion}

\subsection{The Progressive Value of Multi-Source Retrieval}

Our results demonstrate a clear progressive improvement from each retrieval modality across both datasets, supporting the multi-source ``Second Brain'' design principle. Each retrieval layer captures a different aspect of the disease-symptom relationship:

\begin{itemize}
    \item \textbf{Dense embeddings} capture broad semantic similarity---diseases with symptom profiles that ``sound like'' the query.
    \item \textbf{BM25+PPMI} captures statistical co-occurrence patterns---diseases that frequently co-occur with specific symptoms in clinical data, even when the semantic similarity is not obvious.
    \item \textbf{Knowledge graph} captures structural relationships---the ontological ``wiring'' of which symptoms are definitionally associated with which diseases, weighted by multiple evidence sources.
\end{itemize}

This complementarity mirrors how human clinicians reason: they combine textbook knowledge (ontological), clinical experience (co-occurrence), and pattern matching (semantic similarity). The ``Second Brain'' framework formalizes this multi-modal reasoning for AI systems.

The S2D results further illuminate the interplay between retrieval modalities. On S2D, BM25+PPMI underperforms Dense RAG (unlike DDXPlus), indicating that the relative effectiveness of sparse vs.\ dense retrieval depends on the input format: structured symptom codes favor lexical matching, while free-text descriptions favor semantic embeddings. The Multi-Source KG compensates for this by integrating both signals, achieving the best performance regardless of input format.

\subsection{Model-Agnostic Retrieval: RAG as an Equalizer}

The GPT-4o ablation reveals a surprising and practically important finding: RAG augmentation acts as an \emph{equalizer} across model capabilities. The 9.0-percentage-point gap between GPT-4o-mini and GPT-4o in the LLM-Only condition (32.5\% vs.\ 23.5\%) vanishes entirely with Multi-Source KG augmentation (both at 95.5\% Top-1, 98.5\% Top-5). This suggests that for well-defined diagnostic tasks:

\begin{enumerate}
    \item The base model's parametric medical knowledge matters less than the quality of retrieved evidence.
    \item Smaller, cheaper models can match larger models when supported by appropriate retrieval infrastructure.
    \item The retrieval framework provides a ``floor'' of performance that is largely independent of model size.
\end{enumerate}

This finding has significant practical implications: organizations can deploy smaller, more cost-effective models without sacrificing diagnostic accuracy, provided they invest in high-quality knowledge bases and retrieval pipelines.

\subsection{Domain-Specific Knowledge Graphs}

The strong S2D results (90.3\% Top-1, 98.8\% Top-5) demonstrate that our pipeline generalizes to new medical domains when a domain-specific knowledge graph is constructed from available training data. The KG construction process---computing co-occurrence statistics and PPMI scores from the S2D training split (70\%)---is fully automatic and requires no expert curation, making it practical for deployment across diverse medical specialties.

However, the performance gap between DDXPlus (95.5\% Top-1) and S2D (90.3\% Top-1) reflects genuine differences in task difficulty: S2D uses free-text symptom descriptions with high linguistic variability, fewer training examples per disease (approximately 31 vs.\ 204 for DDXPlus), and a less structured knowledge representation. These factors limit the quality of the automatically constructed KG, suggesting that hybrid approaches combining automatic construction with expert curation could further improve performance.

\subsection{TF-IDF and PPMI as Clinical Significance Proxies}

Our use of TF-IDF and PPMI as proxies for clinical significance of symptom-disease associations is effective but imperfect. These statistical measures capture population-level patterns but may not reflect:

\begin{itemize}
    \item \textbf{Pathognomonic symptoms:} Highly specific indicators (e.g., Koplik spots for measles) that may be rare in training data but clinically definitive.
    \item \textbf{Severity weighting:} The clinical urgency associated with certain symptom-disease combinations.
    \item \textbf{Temporal patterns:} The diagnostic significance of symptom progression over time.
\end{itemize}

Integrating expert-curated specificity weights from clinical knowledge bases (e.g., DXplain, Isabel) could address these limitations.

\subsection{Practical Implications for Clinical Decision Support}

The near-perfect performance on both benchmarks (98.5\% and 98.8\% Top-5, zero hallucination) suggests that multi-source RAG systems could serve as reliable clinical decision support tools within well-defined domains. Key deployment considerations include:

\begin{itemize}
    \item \textbf{Knowledge base maintenance:} Regular updates as new diseases and symptom associations are identified.
    \item \textbf{Scope transparency:} Clear communication to clinicians about which diseases and symptoms the system covers.
    \item \textbf{Domain adaptation:} Automatic KG construction from local clinical data to adapt the system to institution-specific disease distributions.
    \item \textbf{Regulatory compliance:} Adherence to medical device regulations for clinical decision support software.
\end{itemize}

%% ============================================================
\section{Limitations}
\label{sec:limitations}

We acknowledge several limitations of this work:

\textbf{Dataset scope.} DDXPlus covers 49 diseases and S2D covers 24 diseases, both subsets of real clinical pathology. DDXPlus cases are synthetically generated from probabilistic models and may not capture the full complexity of real clinical presentations, including comorbidities, atypical presentations, and incomplete information.

\textbf{Sample size.} Our DDXPlus evaluation uses $n{=}200$ stratified samples and S2D uses $n{=}320$ test cases. While sufficient to demonstrate large effect sizes, larger evaluation samples would provide tighter confidence intervals.

\textbf{Two-model evaluation.} We evaluate GPT-4o-mini and GPT-4o. The interaction between retrieval augmentation and model capability may differ for open-source medical models (e.g., Meditron~\cite{chen2023meditron}, BioMistral) or other architectures.

\textbf{Static knowledge graph.} Our KGs are constructed from training data and do not incorporate external medical knowledge bases (UMLS, SNOMED-CT) or adapt to new information over time.

\textbf{Equal fusion weights.} We use equal weights ($\alpha{=}\beta{=}\gamma{=}1/3$) for KG edge types. Learned or tuned weights could potentially improve performance, though the near-perfect results leave limited room for improvement.

%% ============================================================
\section{Conclusion}
\label{sec:conclusion}

We presented a multi-source RAG framework that serves as a ``Second Brain'' for LLM-based medical diagnosis. By progressively integrating dense vector retrieval, BM25 with PPMI co-occurrence statistics, and a multi-source knowledge graph, our system achieves near-perfect diagnostic accuracy with zero hallucinations across two medical benchmarks: 95.5\% Top-1 and 98.5\% Top-5 on DDXPlus, and 90.3\% Top-1 and 98.8\% Top-5 on Symptom2Disease. This compares to unaugmented LLM baselines of 32.5\% and 32.5\% Top-1 accuracy with 59.8\% and 75.9\% hallucination rates, respectively.

Each retrieval layer contributes complementary information in a clear progressive pattern: dense retrieval provides the largest accuracy gain, sparse retrieval with PPMI captures statistical patterns missed by embeddings, and the knowledge graph optimizes ranking quality while eliminating residual hallucinations. A model ablation with GPT-4o demonstrates that our retrieval framework compensates for variation in base model capability, with both GPT-4o and GPT-4o-mini converging to identical performance under full augmentation.

The cross-domain evaluation on Symptom2Disease demonstrates that our pipeline generalizes to new medical domains when domain-specific knowledge graphs are constructed from available training data---an automatic process requiring no expert curation. The Multi-Source KG provides its largest relative improvement precisely when individual retrieval modalities show mixed performance, validating the multi-source design principle.

The multi-source RAG paradigm offers a principled path toward reliable LLM-based clinical decision support---not as an autonomous diagnostician, but as a grounded, transparent ``Second Brain'' that augments clinical reasoning with structured, multi-modal medical knowledge.

%% ============================================================
\section*{CRediT Authorship Contribution Statement}

\textbf{First Author:} Conceptualization, Methodology, Software, Validation, Writing -- Original Draft.
\textbf{Second Author:} Methodology, Formal Analysis, Writing -- Review \& Editing.
\textbf{Third Author:} Supervision, Writing -- Review \& Editing.

\section*{Declaration of Competing Interest}

The authors declare that they have no known competing financial interests or personal relationships that could have appeared to influence the work reported in this paper.

\section*{Data Availability}

The DDXPlus dataset is publicly available at \url{https://github.com/mila-iqia/ddxplus}. The Symptom2Disease dataset is available at \url{https://huggingface.co/datasets/gretelai/symptom_to_diagnosis}. Code and knowledge graph construction scripts will be released upon publication.

%% ============================================================
\bibliographystyle{elsarticle-num}

\begin{thebibliography}{25}

\bibitem{brown2020language}
T.~Brown, B.~Mann, N.~Ryder, et~al., Language models are few-shot learners, in: Advances in Neural Information Processing Systems, vol.~33, 2020, pp.~1877--1901.

\bibitem{openai2023gpt4}
OpenAI, GPT-4 technical report, arXiv preprint arXiv:2303.08774, 2023.

\bibitem{ji2023survey}
Z.~Ji, N.~Lee, R.~Frieske, et~al., Survey of hallucination in natural language generation, ACM Computing Surveys 55~(12) (2023) 1--38.

\bibitem{singhal2023large}
K.~Singhal, S.~Azizi, T.~Tu, et~al., Large language models encode clinical knowledge, Nature 620~(7972) (2023) 172--180.

\bibitem{lewis2020retrieval}
P.~Lewis, E.~Perez, A.~Piktus, et~al., Retrieval-augmented generation for knowledge-intensive NLP tasks, in: Advances in Neural Information Processing Systems, vol.~33, 2020, pp.~9459--9474.

\bibitem{gao2023retrieval}
Y.~Gao, Y.~Xiong, A.~Gupta, et~al., Retrieval-augmented generation for large language models: A survey, arXiv preprint arXiv:2312.10997, 2023.

\bibitem{forte2022building}
T.~Forte, Building a Second Brain: A Proven Method to Organize Your Digital Life and Unlock Your Creative Potential, Atria Books, 2022.

\bibitem{robertson2009probabilistic}
S.~Robertson, H.~Zaragoza, The probabilistic relevance framework: BM25 and beyond, Foundations and Trends in Information Retrieval 3~(4) (2009) 333--389.

\bibitem{fansi2022ddxplus}
A.~Fansi~Tchango, R.~Camara, I.~Kreitchmann, et~al., DDXPlus: A new dataset for automatic medical diagnosis, in: Advances in Neural Information Processing Systems, vol.~35, 2022.

\bibitem{gretelai_s2d}
Gretel.ai, Symptom to diagnosis dataset, Hugging Face Datasets, 2023. URL: \url{https://huggingface.co/datasets/gretelai/symptom_to_diagnosis}.

\bibitem{santos2022knowledge}
A.~Santos, A.~Coelho, et~al., A knowledge graph to interpret clinical proteomics data, Nature Biotechnology 40 (2022) 692--702.

\bibitem{chandak2023building}
P.~Chandak, K.~Huang, M.~Zitnik, Building a knowledge graph to enable precision medicine, Scientific Data 10 (2023) 67.

\bibitem{bodenreider2004unified}
O.~Bodenreider, The Unified Medical Language System (UMLS): Integrating biomedical terminology, Nucleic Acids Research 32 (2004) D267--D270.

\bibitem{pan2024unifying}
S.~Pan, L.~Luo, Y.~Wang, et~al., Unifying large language models and knowledge graphs: A roadmap, IEEE Transactions on Knowledge and Data Engineering (2024).

\bibitem{mcduff2023towards}
D.~McDuff, M.~Schaekermann, T.~Tu, et~al., Towards accurate differential diagnosis with large language models, arXiv preprint arXiv:2312.00164, 2023.

\bibitem{karpukhin2020dense}
V.~Karpukhin, B.~O\u{g}uz, S.~Min, et~al., Dense passage retrieval for open-domain question answering, in: EMNLP, 2020, pp.~6769--6781.

\bibitem{izacard2021leveraging}
G.~Izacard, E.~Grave, Leveraging passage retrieval with generative models for open domain question answering, in: EACL, 2021, pp.~874--880.

\bibitem{shuster2021retrieval}
K.~Shuster, S.~Poff, M.~Chen, et~al., Retrieval augmentation reduces hallucination in conversation, in: EMNLP Findings, 2021.

\bibitem{thorne2018fever}
J.~Thorne, A.~Vlachos, C.~Christodoulopoulos, A.~Mittal, FEVER: A large-scale dataset for fact extraction and verification, in: NAACL-HLT, 2018, pp.~809--819.

\bibitem{zhao2024retrieval}
P.~Zhao, H.~Zhang, Q.~Yu, et~al., Retrieval-augmented generation for AI-generated content: A survey, arXiv preprint arXiv:2402.19473, 2024.

\bibitem{umapathi2023med}
L.~K.~Umapathi, A.~Pal, M.~Sankarasubbu, Med-HALT: Medical domain hallucination test for large language models, in: CoNLL, 2023.

\bibitem{lu2022neurologic}
X.~Lu, S.~Welleck, P.~West, et~al., NeuroLogic A*esque decoding: Constrained text generation with lookahead heuristics, in: NAACL, 2022.

\bibitem{manakul2023selfcheckgpt}
P.~Manakul, A.~Liusie, M.~J.~F.~Gales, SelfCheckGPT: Zero-resource black-box hallucination detection for generative large language models, in: EMNLP, 2023.

\bibitem{cormack2009reciprocal}
G.~V.~Cormack, C.~L.~A.~Clarke, S.~B\"uttcher, Reciprocal rank fusion outperforms condorcet and individual rank learning methods, in: SIGIR, 2009, pp.~758--759.

\bibitem{he2024gretriever}
X.~He, Y.~Tian, Y.~Sun, et~al., G-Retriever: Retrieval-augmented generation for textual graph understanding and question answering, arXiv preprint arXiv:2402.07630, 2024.

\bibitem{yasunaga2022deep}
M.~Yasunaga, A.~Bosselut, H.~Ren, et~al., Deep bidirectional language-knowledge graph pretraining, in: NeurIPS, 2022.

\bibitem{chen2023meditron}
Z.~Chen, A.~Cano, et~al., Meditron-70B: Scaling medical pretraining for large language models, arXiv preprint arXiv:2311.16079, 2023.

\end{thebibliography}

\end{document}
